\chapter{Conclusion and Future Work}

% ─────────────────────────────────────────────────────────────────────────────
\section{Summary of the Work}

This thesis has presented the design, development, and evaluation of a
cloud-native H5P interactive video content creation platform for university
instructors, with a dual focus on user-centred workflow design and
AI-assisted question generation.

The core technical contribution is a full-stack web application comprising
a React~18 + TypeScript single-page application, an Express.js REST API,
a SQLite database managed via Sequelize ORM, and the Lumieducation H5P
server library. Layered on top of this foundation is an AI pipeline that
extracts timestamped transcripts from videos (via YouTube captions, WebVTT
uploads, or Whisper audio transcription), submits them to a large language
model (Anthropic Claude or a locally hosted Ollama instance), and returns
a Zod-validated list of H5P interaction suggestions classified by Bloom's
revised taxonomy. A Server-Sent Events streaming interface delivers token-level
AI progress updates to the browser, providing immediate feedback during the
8--12 second generation latency.

The platform was evaluated against the dominant WordPress-based H5P workflow
in a comparative usability study with ten participants. The new platform
achieved a 95\% task completion rate (versus 30\% for the baseline), a median
task time of 3.2 minutes (versus 45 minutes), a mean error count of 0.9
(versus 8.3), a mean SUS score of 82.5 (versus 38.0), and a mean
satisfaction rating of 4.8/5 (versus 2.1/5). These results represent
substantial improvements across every metric and exceed all non-functional
requirements set at the outset of the project.

% ─────────────────────────────────────────────────────────────────────────────
\section{Answers to Research Questions}

\textbf{RQ1: What are the specific usability barriers in current H5P
authoring workflows?}

The primary barriers identified were multi-screen navigation fragmentation,
the requirement to understand CMS concepts irrelevant to pedagogy, the
absence of real-time preview, and the absence of any mechanism for obtaining
question suggestions. All four were directly addressed in the new platform.

\textbf{RQ2: How can user-centred design reduce extraneous cognitive load?}

Co-locating the video player, timestamp capture, question editor, and AI
suggestion panel on a single screen eliminated the split-attention effect
identified in the cognitive load literature. The integrated format reduced
the number of required interface interactions from 47 (WordPress) to
approximately 8 for a three-question interactive video.

\textbf{RQ3: What architecture best supports AI-assisted H5P authoring?}

The dual-provider AI architecture (Ollama as first choice, Claude as
fallback) combined with Zod-based structured output validation provides a
reliable, cost-flexible pipeline. SSE streaming is essential for managing
perceived latency during the 8--12 second generation time. The Lumieducation
H5P server library provides a robust, specification-compliant H5P layer
without CMS dependency.

\textbf{RQ4: To what extent does the platform outperform the baseline?}

The platform outperforms the WordPress baseline on every measured dimension:
task completion rate by 65 percentage points, time-on-task by 93\%, error
rate by 89\%, SUS score by 44.5 points, and satisfaction by 2.7 points on
a five-point scale. These improvements are practically significant and
suggest that the platform is ready for pilot deployment at VNU-HCM.

% ─────────────────────────────────────────────────────────────────────────────
\section{Contributions}

This thesis makes the following contributions:

\begin{enumerate}
  \item \textbf{A production-deployable H5P authoring platform.} The
  platform provides a complete, open-source, CMS-free H5P authoring
  environment suitable for institutional deployment. It is the first
  publicly documented system to combine the Lumieducation H5P server
  library with a modern React SPA and a JWT-secured REST API.

  \item \textbf{An AI pipeline for H5P question generation.} The
  transcript-to-H5P suggestion pipeline is the first documented integration
  of a production-grade LLM (Claude/Ollama) with structured H5P content
  type generation, Bloom's taxonomy classification, and runtime schema
  validation. The pipeline design is modular and can be adapted to other
  educational content formats.

  \item \textbf{Prompt engineering patterns for structured H5P output.}
  The system prompt, output automater pattern, and Zod validation strategy
  developed in this project constitute a reusable template for generating
  structured educational content with LLMs. The specific challenge of
  ensuring that LLM-generated JSON conforms to H5P library schemas is
  documented and solved in a reproducible way.

  \item \textbf{Empirical evidence for user-centred design in EdTech.}
  The comparative evaluation provides quantitative evidence that a
  purpose-built, user-centred platform can dramatically outperform a
  general-purpose CMS-based workflow for educational authoring tasks,
  contributing to the evidence base for user-centred design in educational
  technology.

  \item \textbf{An alignment with Vietnam's digital education strategy.}
  The platform supports VNU-HCM's digital transformation goals by providing
  an open-source, locally deployable tool that does not depend on commercial
  licensing and is compatible with the existing MOOC infrastructure.
\end{enumerate}

% ─────────────────────────────────────────────────────────────────────────────
\section{Limitations}

Several limitations constrain the interpretation of the findings.

\textbf{Evaluation sample size.} Ten participants is sufficient for a
heuristic usability study but insufficient for statistical hypothesis
testing~\cite{nielsen1994}. The reported effect sizes are large enough to
be practically interpretable, but formal significance testing would require
a larger and more representative sample including actual university faculty.

\textbf{Proxy participants.} IT students were used as educator proxies.
While standard practice for early-stage evaluation, this introduces a gap
in ecological validity. Faculty members may exhibit different error patterns
and different responses to the AI suggestion interface.

\textbf{Short-term evaluation.} The study measured a single-session task.
Long-term adoption---whether instructors continue to use the platform after
initial exposure, whether content quality improves over time, and whether
the AI pipeline remains relevant as video content evolves---was not assessed.

\textbf{Single-institution context.} The requirements elicitation and
participant pool were anchored at HCMIU. Generalisability to other
Vietnamese universities, or to non-Vietnamese educational contexts, requires
validation.

\textbf{AI suggestion quality variance.} Suggestion quality depends on
transcript quality, which degrades with noisy audio or heavy technical
vocabulary. The current prompt strategy has not been optimised for
non-English videos; Vietnamese-language lecture transcription and question
generation are important future directions.

% ─────────────────────────────────────────────────────────────────────────────
\section{Future Work}

The following extensions are identified as high-priority directions for
future research and development.

\textbf{Faculty evaluation study.} A longitudinal usability study with
actual university faculty, measuring adoption over a full semester, would
provide ecologically valid evidence for the platform's educational impact.

\textbf{Vietnamese language support.} Extending the AI pipeline to support
Vietnamese-language transcripts and generate questions in Vietnamese would
make the platform accessible to a much larger fraction of VNU-HCM faculty.
This requires both Whisper fine-tuning on Vietnamese academic audio and
prompt adaptation for Vietnamese H5P content generation.

\textbf{Learning analytics integration.} Connecting the platform's export
events with the VNU-HCM MOOC LRS (Learning Record Store) would enable
instructors to see how students interact with their interactive videos---
which questions cause the most replays, which timestamps trigger the most
errors, and how interaction patterns correlate with assessment outcomes.
This integration would directly address the learning analytics gap
identified by Siemens and Long~\cite{siemens2011}.

\textbf{Collaborative authoring.} Adding real-time collaborative editing
(multiple instructors working on the same video simultaneously) via
Socket.io's room-based messaging would enable team-based content production.
The Socket.io dependency is already included in the backend and the
architecture supports this extension.

\textbf{AI fine-tuning on educational content.} The current pipeline uses
general-purpose LLMs with prompt engineering. Fine-tuning a smaller model
(e.g., a Llama variant) on a corpus of validated H5P question--transcript
pairs from the VNU-HCM MOOC platform would improve suggestion quality and
reduce generation latency, particularly for the Ollama local inference path.

\textbf{Asynchronous video processing queue.} The current video upload
pipeline performs FFmpeg thumbnail generation synchronously in the request
handler, causing 220--410\,ms response latencies for the upload endpoint.
Implementing an asynchronous job queue (e.g., using Bull or BullMQ) would
reduce upload endpoint latency and improve system responsiveness under load.

\textbf{Full WCAG~2.1 AA compliance audit.} The current platform includes
basic accessibility features (keyboard navigation, ARIA labels on interactive
elements) but has not undergone a formal WCAG~2.1 AA compliance audit.
Achieving full accessibility compliance is essential before the platform can
be offered as an institutional standard tool.
